% Part: propositional-logic
% Chapter: syntax-and-semantics
% Section: valuations-sat

\documentclass[../../../include/open-logic-section]{subfiles}

\begin{document}

\olfileid{pl}{syn}{val}

\olsection{\usetoken{P}{valuation} ও পরিতৃপ্তি}

\begin{defn}[!!^{valuation}s]
“সত্য” ও “মিথ্যা”—এই দুই সত্যমানের সেট হোক $\{\True, \False\}$।
$\Lang{L_0}$-এর জন্য একটি \emph{!!{valuation}} হলো এমন একটি
অপেক্ষক~$\pAssign{v}$, যা ভাষাটির !!{propositional variable}s-কে
$\True$ অথবা $\False$ আরোপ করে; অর্থাৎ $\pAssign{v} \colon
\PVar \to \{\True, \False \}$।
\end{defn}

\begin{defn}
  কোনো !!a{valuation}~$\pAssign{v}$ দেওয়া থাকলে, মূল্যায়ন অপেক্ষক
  $\pValue{v} \colon \Frm[L_0] \to \{\True, \False \}$-কে নিচের মতো
  আরোহীভাবে সংজ্ঞায়িত করি:
  \begin{align*}
    \iftag{prvFalse}{\pValue{v}(\lfalse) & = \False; \\}{}
    \iftag{prvTrue}{\pValue{v}(\ltrue) & = \True; \\}{}
    \pValue{v}(\Obj p_n) & = \pAssign{v}(\Obj p_n); \\
    \iftag{prvNot}{\pValue{v}(\lnot !A) & = \begin{cases}
        \True & \text{যদি } \pValue{v}(!A) = \False;\\
        \False & \text{অন্যথায়।}
      \end{cases} \\ }{}
    \iftag{prvAnd}{\pValue{v}(!A \land !B) & = \begin{cases}
        \True &
        \text{যদি $\pValue{v}(!A) = \True$ এবং $\pValue{v}(!B) = \True$;}\\
        \False &
        \text{যদি $\pValue{v}(!A) = \False$ অথবা $\pValue{v}(!B) = \False$}.
      \end{cases}\\}{}
    \iftag{prvOr}{\pValue{v}(!A \lor !B) & = \begin{cases}
        \True &
        \text{যদি $\pValue{v}(!A) = \True$ অথবা $\pValue{v}(!B) = \True$;}\\
        \False &
        \text{যদি $\pValue{v}(!A) = \False$ এবং $\pValue{v}(!B) = \False$}.
      \end{cases}\\}{}
    \iftag{prvIf}{\pValue{v}(!A \lif !B) & = \begin{cases}
        \True &
        \text{যদি $\pValue{v}(!A) = \False$ অথবা $\pValue{v}(!B) = \True$;}\\
        \False &
        \text{যদি $\pValue{v}(!A) = \True$ এবং $\pValue{v}(!B) = \False$}.
      \end{cases}\\}{}
    \iftag{prvIff}{\pValue{v}(!A \liff !B) & = \begin{cases}
        \True &
        \text{যদি $\pValue{v}(!A) = \pValue{v}(!B)$;}\\
        \False &
        \text{যদি $\pValue{v}(!A) \neq \pValue{v}(!B)$}.
      \end{cases}}{}
\end{align*}
\end{defn}

\begin{explain}
বিধিগুলি নিচের সত্যসারণিগুলির সঙ্গে মেলে:
\begin{center}
\iftag{prvNot}{
  \begin{tabular}{|c||c|} \hline
    $!A$ & $ \lnot !A$ \\
    \hline \hline
    $\True$ & $\False$ \\
    $\False$ & $\True$ \\
    \hline
  \end{tabular}
}{}
\iftag{prvAnd}{
  \begin{tabular}{|cc||c|} \hline
    $!A$ & $!B$ & $!A \land !B$ \\
    \hline \hline
    $\True$ & $\True$ & $\True$ \\
    $\True$ & $\False$ & $\False$ \\
    $\False$ & $\True$ & $\False$ \\
    $\False$ & $\False$ & $\False$ \\
    \hline
  \end{tabular}
}{}
\iftag{prvOr}{
  \begin{tabular}{|cc||c|} \hline
    $!A$ & $!B$ & $!A \lor !B$ \\
    \hline \hline
    $\True$ & $\True$ & $\True$ \\
    $\True$ & $\False$ & $\True$ \\
    $\False$ & $\True$ & $\True$ \\
    $\False$ & $\False$ & $\False$ \\
    \hline
  \end{tabular}
}{}%
\iftag{notprvNot,notprvAnd,notprvOr,notprvIf}{}{\\[1em]}
\iftag{prvIf}{
  \begin{tabular}{|cc||c|} \hline
    $!A$ & $!B$ & $!A \lif !B$ \\
    \hline \hline
    $\True$ & $\True$ & $\True$ \\
    $\True$ & $\False$ & $\False$ \\
    $\False$ & $\True$ & $\True$ \\
    $\False$ & $\False$ & $\True$ \\
    \hline
  \end{tabular}
}{}
\iftag{prvIff}{
  \begin{tabular}{|cc||c|} \hline
    $!A$ & $!B$ & $!A \liff !B$ \\
    \hline \hline
    $\True$ & $\True$ & $\True$ \\
    $\True$ & $\False$ & $\False$ \\
    $\False$ & $\True$ & $\False$ \\
    $\False$ & $\False$ & $\True$ \\
    \hline
  \end{tabular}
}{}
\end{center}
\end{explain}

\begin{prob}
$\Lang{L_0}$-এ একটি ত্রিস্থানী সংযোজক $\diamondsuit$ যোগ করার কথা ভাবো,
যার মূল্যায়ন নিচে দেওয়া:
\begin{gather*}
  \pValue{v}(\diamondsuit( !A , !B , !C ) ) = \begin{cases}
    \pValue{v}( !B ) &
    \text{যদি $\pValue{v}(!A) = \True$;}\\
    \pValue{v}( !C ) &
    \text{যদি $\pValue{v}(!A) = \False $}.
  \end{cases}
\end{gather*}
এই সংযোজকের সত্যসারণি লেখো।
\end{prob}

\begin{thm}[স্থানীয় নির্ধারণ]
\ollabel{thm:LocalDetermination} ধরি $\pAssign{v_1}$ ও
$\pAssign{v_2}$ এমন দুটি !!{valuation}s, যারা $!A$-তে উপস্থিত
!!{propositional
variable}s-গুলির উপর একই মান দেয়; অর্থাৎ যখনই
$\Obj p_n$ !!{formula}~$!A$-টিতে উপস্থিত, তখন
$\pAssign{v_1}(\Obj p_n) = \pAssign{v_2}(\Obj p_n)$। তাহলে
$\pValue{v_1}$ ও $\pValue{v_2}$-ও~$!A$-এর উপর একই মান দেয়; অর্থাৎ
$\pValue{v_1}(!A) = \pValue{v_2}(!A)$।
\end{thm}

\begin{proof}
$!A$-এর উপর আরোহ প্রয়োগ করি।
\end{proof}

\begin{defn}[পরিতৃপ্তি]
\ollabel{defn:satisfaction} কোনো
  \emph{!!a{valuation}~$\pAssign{v}$ দ্বারা !!a{formula}~$!A$-এর
  পরিতৃপ্তি}, $\pSat{v}{!A}$, ধারণাটি নিচের মতো আরোহীভাবে সংজ্ঞায়িত
  করা যায়। (“$\pSat{v}{!A}$ নয়” বোঝাতে আমরা $\pSat/{v}{!A}$ লিখি।)
\begin{enumerate}
\tagitem{prvFalse}{%
  \indcase{!A}{\lfalse}{$\pSat/{v}{\indfrm}$।}}{}

\tagitem{prvTrue}{%
  \indcase{!A}{\ltrue}{$\pSat{v}{\indfrm}$।}}{}

\item \indcase{!A}{\Obj p_i}{$\pSat{v}{\indfrm}$ তখনই,
  যখন $\pAssign{v}(\Obj p_i) = \True$।}

\tagitem{prvNot}{%
  \indcase{!A}{\lnot !B}{$\pSat{v}{\indfrm}$ তখনই, যখন
    $\pSat/{v}{!B}$।}}{}

\tagitem{prvAnd}{%
  \indcase{!A}{(!B \land !C)}{$\pSat{v}{\indfrm}$ তখনই, যখন
    $\pSat{v}{!B}$ এবং $\pSat{v}{!C}$।}}{}

\tagitem{prvOr}{%
  \indcase{!A}{(!B \lor !C)}{$\pSat{v}{\indfrm}$ তখনই, যখন
    $\pSat{v}{!B}$ অথবা $\pSat{v}{!C}$ (অথবা উভয়ই)।}}{}

\tagitem{prvIf}{%
  \indcase{!A}{(!B \lif !C)}{$\pSat{v}{\indfrm}$ তখনই, যখন
    $\pSat/{v}{!B}$ অথবা $\pSat{v}{!C}$ (অথবা উভয়ই)।}}{}

\tagitem{prvIff}{%
  \indcase{!A}{(!B \liff !C)}{$\pSat{v}{\indfrm}$ তখনই, যখন হয়
    $\pSat{v}{!B}$ ও $\pSat{v}{!C}$ উভয়ই সত্য, নয়তো
    $\pSat{v}{!B}$ ও $\pSat{v}{!C}$ কোনোটিই সত্য নয়।}}{}
\end{enumerate}
$\Gamma$ যদি !!{formula}s-এর একটি সেট হয়, তবে $\pSat{v}{\Gamma}$
তখনই, যখন প্রতিটি~$!A \in \Gamma$-এর জন্য $\pSat{v}{!A}$।
\end{defn}

\begin{prop}\ollabel{prop:sat-value}
  $\pSat{v}{!A}$ তখনই, যখন $\pValue{v}(!A) = \True$।
\end{prop}

\begin{proof}
  $!A$-এর উপর আরোহ প্রয়োগ করি।
\end{proof}

\begin{prob}
  \olref[pl][syn][val]{prop:sat-value} প্রমাণ করো।
\end{prob}

\end{document}
